\chapter{Onglet Paramétrage XFoil}
\label{chap:xfoil}

L'onglet \menu{Paramétrage XFoil} regroupe l'ensemble des paramètres
de simulation transmis au solveur. Il est structuré en zones de
saisie (Reynolds, Alpha, Paramètres XFoil), une rangée de boutons
de lancement et un cadre \menu{Diagnostic}.

\screenshot{09_tab_xfoil.png}{1.0}{%
    Onglet \menu{Paramétrage XFoil} : zones de saisie, boutons de
    lancement (dont \og Flap seul \fg{}) et cadre \og Diagnostic \fg{}.}

\begin{noteinfo}
Valeurs par défaut au démarrage : Reynolds de 1\,000\,000 à
2\,000\,000 (pas 1\,000\,000), incidence de $-5^\circ$ à $+15^\circ$
(pas $0{,}5^\circ$), transition forcée XTR~$=0{,}2$, 100~itérations,
balayage \menu{ASEQ}.
\end{noteinfo}

\begin{noteinfo}
Ces paramètres alimentent le \textbf{solveur sélectionné} dans
\menu{Options $\rightarrow$ Solveur} (XFoil ou FlexFoil, voir
chapitre~\ref{chap:interface}). Les deux moteurs interprètent les mêmes
réglages ; certaines options propres à l'exécutable XFoil (balayage
\menu{ASEQ}, réinitialisation en cas d'échec) n'ont pas d'effet sur
FlexFoil, qui converge par une méthode native.
\end{noteinfo}

\section{Plage de Reynolds}

La zone \menu{Reynolds} permet de définir une \textbf{plage} de
nombres de Reynolds pour le balayage. Pour chaque Reynolds, une
polaire complète est calculée. Trois champs :

\begin{tabularx}{\linewidth}{l X}
    \toprule
    \textbf{Champ} & \textbf{Description} \\
    \midrule
    Min  & Reynolds minimum du balayage. \\
    Max  & Reynolds maximum du balayage. \\
    Pas  & Incrément entre deux valeurs consécutives. Pour
           ne calculer qu'un seul Reynolds, choisir
           \texttt{Pas $\geq$ Max - Min}. \\
    \bottomrule
\end{tabularx}



\section{Plage d'incidence (Alpha)}

La zone \menu{Alpha (deg)} définit le balayage en angle d'incidence
$\alpha$ pour chaque polaire :

\begin{tabularx}{\linewidth}{l X}
    \toprule
    \textbf{Champ} & \textbf{Description} \\
    \midrule
    Min  & Incidence minimum (en degrés, peut être négative). \\
    Max  & Incidence maximum. \\
    Pas  & Incrément entre deux points de la polaire. Plus
           petit~=~polaire plus fine, mais simulation plus longue.\\
    \bottomrule
\end{tabularx}

\subsection*{Mode \menu{ASEQ}}

La case à cocher \menu{ASEQ} contrôle la \textbf{stratégie de
balayage} en incidence~:

\begin{itemize}
    \item \textbf{Cochée} (par défaut) : XFoil exécute la commande
        \texttt{ASEQ} qui balaye automatiquement la plage en
        partant du minimum. C'est le mode \emph{rapide} historique.
    \item \textbf{Décochée} : AirfoilTools envoie une succession de
        commandes \texttt{ALFA} individuelles, en partant de
        $\alpha=0$ vers le maximum, puis depuis $-\Delta$ vers le
        minimum (avec une commande \texttt{INIT} entre les deux).
        Ce mode est plus \emph{robuste} à la convergence aux
        angles élevés ou pour des profils difficiles.
\end{itemize}

\begin{astuce}
Si XFoil ne converge pas pour certains angles élevés, désactivez
\menu{ASEQ} pour tenter le balayage bidirectionnel. Cela
améliore souvent la convergence en initialisant chaque branche
à partir de la solution à $\alpha=0$.
\end{astuce}

\subsection*{Option \menu{Réinit. si échec}}

Lorsqu'un point d'incidence ne converge pas (message
\texttt{VISCAL: Convergence failed}), le balayage en marche peut
\textbf{abandonner toute la suite} de la branche (effet de cascade).
La case \menu{Réinit. si échec} ajoute une \textbf{seconde passe}
qui rejoue chaque incidence depuis un état propre (commande
\texttt{INIT} avant chaque \texttt{ALFA}), afin de récupérer les
points manqués sans dégrader le balayage en marche (le
postprocesseur conserve la solution obtenue en marche lorsqu'elle
existe).

\begin{attention}
L'option \menu{Réinit. si échec} implique le mode \texttt{ALFA}
individuel (incompatible avec \menu{ASEQ}) et \textbf{double
environ} le temps de calcul. Pensez à augmenter le \menu{Timeout}
en conséquence. Elle est \textbf{décochée par défaut}.
\end{attention}

\section{Paramètres XFoil avancés}

La zone \menu{Paramètres XFoil} regroupe les options du solveur
proprement dit.

\subsection{Viscosité}

\begin{itemize}
    \item Cochée (\menu{Avec}) : analyse \textbf{visqueuse}
        (commande \texttt{VISC} de XFoil), avec couche limite.
        \textbf{Recommandé pour toute polaire réaliste.}
    \item Décochée : analyse non-visqueuse (potentielle, sans
        frottement). Utile uniquement pour des comparaisons
        géométriques rapides.
\end{itemize}

\subsection{NCRIT}

Critère de transition de couche limite par la méthode $e^N$
(amplification des ondes de Tollmien-Schlichting).

\begin{tabularx}{\linewidth}{c X}
    \toprule
    \textbf{NCRIT} & \textbf{Type d'écoulement} \\
    \midrule
    9       & Écoulement libre standard (valeur par défaut). \\
    4--7    & Écoulement très turbulent (souffleries fermées,
              forte turbulence ambiante). \\
    12--14  & Écoulement très calme (vol planeur en air
              parfaitement laminaire). \\
    \bottomrule
\end{tabularx}

\subsection{XTR Top et XTR Bot}

Position de \textbf{transition forcée} de la couche limite,
respectivement sur l'extrados (\emph{top}) et l'intrados
(\emph{bottom}), exprimée en fraction de corde ($x/c$).

\begin{itemize}
    \item Valeur 1{,}0 : transition libre, calculée par le critère
        NCRIT. C'est le mode standard.
    \item Valeur $<1{,}0$ : transition forcée à la position indiquée
        (simule un \emph{trip} turbulent : turbulateur, rugosité,
        impact d'insectes\dots).
\end{itemize}

\begin{important}
La valeur minimale est \textbf{0{,}05}. Forcer la transition trop
près du bord d'attaque (par exemple 0{,}01) \textbf{déstabilise la
couche limite de XFoil} : la routine de transition diverge
(\texttt{NaN}) et entre dans une boucle quasi infinie, jusqu'à
expiration du \menu{Timeout}. La valeur par défaut est
\textbf{0{,}2}.
\end{important}

\subsection{Repanel et NPanel}

\begin{itemize}
    \item \menu{Repanel} : si cochée, XFoil exécute la commande
        \texttt{PANE} qui ré-échantillonne le profil avec un
        critère adaptatif basé sur la courbure. Recommandé pour
        les profils à échantillonnage non-uniforme ou pour des
        profils de très haute résolution.
    \item \menu{NPanel} : nombre de panneaux après ré-échantillonnage
        (par défaut 200). Plus = plus précis mais plus lent.
\end{itemize}

\subsection{Timeout}

Durée maximale d'exécution de XFoil pour le calcul d'une polaire,
en secondes (par défaut 30~s). Au-delà, le processus est tué.
Cela évite à l'application de rester bloquée si XFoil entre dans
une boucle de non-convergence.

\begin{attention}
Augmenter le timeout (par exemple à 120~s) si vous balayez une
plage très large d'incidences ($\Delta \alpha > 30^\circ$), si vous
calculez plusieurs Reynolds, si l'option \menu{Réinit. si échec}
est active (temps doublé), ou si vos profils convergent lentement.
\end{attention}

\subsection{Itérations}

Nombre maximum d'itérations de Newton par point (commande
\texttt{ITER} de XFoil), par défaut 100. \textbf{Augmenter} cette
valeur (200--300) lorsqu'un point précis ne converge pas
(\texttt{VISCAL: Convergence failed}), situation typique près du
décrochage ou avec un coude de volet.

\section{Boutons de lancement}

Quatre boutons en bas de la zone de paramètres permettent de lancer
la simulation :

\begin{itemize}
    \item \bouton{Lancer tout} : simule tous les profils actifs ---
        courant, référence et, si le volet est activé, le profil
        avec volet.
    \item \bouton{Courant seul} : ne simule que le profil courant
        (bleu).
    \item \bouton{Référence seul} : ne simule que le profil de
        référence (rouge).
    \item \bouton{Flap seul} : ne simule que le profil avec volet
        (vert). Nécessite que le volet soit activé dans l'onglet
        \menu{Profils}.
\end{itemize}

Pendant la simulation, les contrôles de l'onglet sont désactivés,
et une \textbf{barre de progression} apparaît dans la barre de
statut : déterminée (un cran par profil terminé) lorsque plusieurs
profils sont simulés, sinon indéterminée (animation « busy »), car un
balayage en incidence s'exécute en un seul calcul sans avancement
intermédiaire. À la fin, l'onglet \menu{Résultats} se sélectionne
automatiquement.

\begin{important}
Si vous relancez une simulation pour \textbf{un seul} des profils,
les résultats antérieurs des autres profils sont \emph{conservés}
(merge). Cela permet par exemple d'itérer rapidement sur le profil
courant tout en gardant la référence fixe pour comparaison.
\end{important}

\section{Diagnostic}
\label{sec:diagnostic}

Le cadre \menu{Diagnostic}, en bas de l'onglet, donne accès aux
fichiers internes de la dernière simulation de chaque profil
(\bouton{Courant}, \bouton{Référence}, \bouton{Volet}). Les
boutons restent inactifs tant qu'aucune simulation n'a été lancée
pour le profil correspondant.

\screenshot{17_dialog_diagnostic.png}{0.85}{%
    Fenêtre de diagnostic : accès au répertoire de travail et aux
    fichiers de la simulation. Le journal \chemin{diagnostic.log}
    récapitule le déroulé du calcul.}

La fenêtre de diagnostic permet de :

\begin{itemize}
    \item parcourir les \textbf{fichiers} de la simulation : journal
        \chemin{diagnostic.log} (résumé du déroulé : étapes,
        avertissements, bilan de convergence), log console XFoil
        (\chemin{xfoil\_alpha.cmd.log}), script de commandes,
        polaires, distributions de pression ;
    \item \textbf{ouvrir le dossier} de travail dans l'explorateur
        de fichiers du système.
\end{itemize}

\begin{astuce}
En cas d'échec ou de résultats incomplets, ouvrez d'abord
\chemin{diagnostic.log} : il indique clairement si le problème
vient d'une non-convergence (\og AUCUN point convergé\dots\fg{}),
d'un \emph{timeout}, ou d'une polaire vide. Le log console XFoil
fournit ensuite le détail incidence par incidence. Le journal est
sauvegardé \textbf{même en cas de timeout}.
\end{astuce}
